\documentclass[11pt,a4paper,oneside]{book}

% ============================================
% PACKAGE IMPORTS
% ============================================

% Encoding and language
\usepackage[utf8]{inputenc}
\usepackage[T1]{fontenc}
\usepackage[english]{babel}

% Page layout
\usepackage[margin=1in]{geometry}

% Line spacing
\usepackage{setspace}
\onehalfspacing

% Graphics
\usepackage{graphicx}
\usepackage{float}
\graphicspath{{images/}}

% Enhanced chapter headings
\usepackage{titlesec}

% Fancy headers and footers
\usepackage{fancyhdr}
\pagestyle{fancy}
\fancyhf{}
\fancyhead[LE,RO]{\thepage}
\fancyhead[RE]{\leftmark}
\fancyhead[LO]{\rightmark}
\renewcommand{\headrulewidth}{0.4pt}

% For quote environment
\usepackage{csquotes}

% Index
\usepackage{makeidx}
\makeindex

% Hyperlinks
\usepackage{hyperref}
\hypersetup{
    colorlinks=true,
    linkcolor=blue,
    filecolor=magenta,
    urlcolor=cyan,
    pdftitle={Book Title},
    pdfauthor={Author Name}
}

% Bibliography (optional)
\usepackage[style=authoryear]{biblatex}
\addbibresource{references.bib}

% ============================================
% CUSTOM COMMANDS
% ============================================

% Epigraph (quote at beginning of chapter)
\newcommand{\chapterquote}[2]{%
    \begin{quotation}
        \em #1
        \begin{flushright}
            --- #2
        \end{flushright}
    \end{quotation}
}

% ============================================
% DOCUMENT INFORMATION
% ============================================

\title{Book Title}
\author{Author Name}
\date{2024}

% ============================================
% DOCUMENT CONTENT
% ============================================

\begin{document}

% ============================================
% TITLE PAGE
% ============================================

\begin{titlepage}
    \centering
    \vspace*{2cm}
    
    % Title
    {\Huge\bfseries Book Title\\[0.5cm]}
    {\Large Subtitle (if applicable)\\[2cm]}
    
    % Author
    {\Large\textsc{Author Name}\\[1.5cm]}
    
    \vfill
    
    % Publisher or additional info
    {\large Publisher Name\\
    City\\
    2024}
\end{titlepage}

% ============================================
% COPYRIGHT PAGE
% ============================================

\newpage
\thispagestyle{empty}
\null
\vfill

\noindent
\textbf{Book Title}\\
\textit{by Author Name}

\vspace{1cm}

\noindent
Copyright \copyright~2024 by Author Name

\vspace{0.5cm}

\noindent
All rights reserved. No part of this publication may be reproduced, distributed, 
or transmitted in any form or by any means, including photocopying, recording, 
or other electronic or mechanical methods, without the prior written permission 
of the publisher, except in the case of brief quotations embodied in critical 
reviews and certain other noncommercial uses permitted by copyright law.

\vspace{1cm}

\noindent
ISBN: 978-X-XXXX-XXXX-X (Paperback)\\
ISBN: 978-X-XXXX-XXXX-X (eBook)

\vspace{1cm}

\noindent
First Edition

\vspace{1cm}

\noindent
Publisher Name\\
Street Address\\
City, State ZIP\\
Country\\
www.publisher-website.com

% ============================================
% FRONT MATTER
% ============================================

\frontmatter

% Dedication
\chapter*{Dedication}
\addcontentsline{toc}{chapter}{Dedication}

\vspace{2cm}
\begin{center}
    \textit{To those who inspired this work}
\end{center}

% Epigraph (optional opening quote)
\cleardoublepage
\thispagestyle{empty}
\vspace*{5cm}
\begin{center}
    \Large\textit{``An inspiring quote that sets the tone for your book.''}\\[1cm]
    \normalsize --- Quote Author
\end{center}

% Table of Contents
\tableofcontents

% List of Figures (optional)
\listoffigures

% List of Tables (optional)
\listoftables

% Preface
\chapter{Preface}

This is the preface where you can explain:
\begin{itemize}
    \item The purpose and scope of the book
    \item Who the intended audience is
    \item How the book is organized
    \item What makes this book unique
    \item Acknowledgments (or put in separate section)
\end{itemize}

% Acknowledgments (optional - can be combined with Preface)
\chapter{Acknowledgments}

Here you thank people who helped make this book possible:
\begin{itemize}
    \item Mentors and advisors
    \item Family and friends
    \item Colleagues and collaborators
    \item Editors and reviewers
    \item Funding sources
\end{itemize}

% ============================================
% MAIN MATTER
% ============================================

\mainmatter

% Part I (optional - for organizing chapters into parts)
\part{First Part Title}

% Chapter 1
\chapter{Introduction}
\label{chap:intro}

% Optional chapter quote
\chapterquote{A relevant quote to introduce the chapter theme.}{Quote Author}

This is the first chapter of your book. It should introduce the main topics 
and themes that will be explored throughout the book.

\section{Background}
\label{sec:background}

Provide context and background information.

\section{Objectives}
\label{sec:objectives}

Explain what the book aims to achieve.

\section{Organization of This Book}
\label{sec:organization}

Briefly describe what each chapter covers.

% Chapter 2
\chapter{Second Chapter Title}
\label{chap:second}

\chapterquote{Another relevant quote for this chapter.}{Quote Author}

Content of the second chapter. You can include:
\begin{itemize}
    \item Text paragraphs
    \item Lists and enumerations
    \item Figures and tables
    \item Code examples
    \item Mathematical equations
\end{itemize}

\section{First Topic}
\label{sec:topic1}

Discuss the first major topic of this chapter.

\subsection{Subtopic}
\label{subsec:subtopic}

Break down complex topics into subsections.

\section{Second Topic}
\label{sec:topic2}

Continue with additional topics.

% Chapter 3
\chapter{Third Chapter Title}
\label{chap:third}

Each chapter should focus on a specific aspect of your book's subject matter.

\section{Figures Example}
\label{sec:figures}

You can include figures like this:

\begin{figure}[h]
    \centering
    % \includegraphics[width=0.8\textwidth]{example-figure}
    \caption{Caption describing the figure.}
    \label{fig:example}
\end{figure}

Reference figures in text: see Figure~\ref{fig:example}.

\section{Tables Example}
\label{sec:tables}

Tables can be created like this:

\begin{table}[h]
    \centering
    \caption{Example table showing data.}
    \label{tab:example}
    \begin{tabular}{lcc}
        \hline
        \textbf{Item} & \textbf{Value 1} & \textbf{Value 2} \\
        \hline
        Sample A & 10 & 20 \\
        Sample B & 15 & 25 \\
        Sample C & 20 & 30 \\
        \hline
    \end{tabular}
\end{table}

% Part II (if using parts)
\part{Second Part Title}

% Chapter 4
\chapter{Fourth Chapter Title}
\label{chap:fourth}

Continue with more chapters as needed for your book structure.

% Chapter 5
\chapter{Fifth Chapter Title}
\label{chap:fifth}

Add as many chapters as your book requires.

% Chapter 6 - Conclusion
\chapter{Conclusion}
\label{chap:conclusion}

Summarize the key points covered in the book and provide final thoughts.

\section{Summary}
\label{sec:summary}

Recap the main themes and findings.

\section{Final Thoughts}
\label{sec:final}

Leave readers with closing remarks and suggestions for further exploration.

% ============================================
% BACK MATTER
% ============================================

\backmatter

% Appendices
\appendix

\chapter{Additional Resources}
\label{app:resources}

Include supplementary material such as:
\begin{itemize}
    \item Further reading recommendations
    \item Useful websites and resources
    \item Tools and software
    \item Extended examples
\end{itemize}

\chapter{Glossary}
\label{app:glossary}

Define technical terms and concepts used throughout the book.

\begin{description}
    \item[Term 1] Definition of term 1
    \item[Term 2] Definition of term 2
    \item[Term 3] Definition of term 3
\end{description}

% Bibliography (optional)
\printbibliography[heading=bibintoc, title={References and Further Reading}]

% Index
\printindex

% About the Author
\chapter*{About the Author}
\addcontentsline{toc}{chapter}{About the Author}

Brief biography of the author, including:
\begin{itemize}
    \item Professional background
    \item Academic credentials
    \item Previous publications
    \item Current work
    \item Contact information or website
\end{itemize}

\end{document}
